\documentclass[11pt]{article}
\usepackage[margin=1in]{geometry}
\usepackage{amsmath,graphicx,hyperref,natbib}
\title{Astronomical CCDs as Long-Baseline Space-Radiation Dosimeters:\\
Cross-Mission Inference from HST, Gaia and Euclid}
\author{Biswajit Jana}
\date{Methods and Data working draft}
\begin{document}
\maketitle
\begin{abstract}
This document records the data and extraction methods of an initial reproducibility
pilot. A calibrated longitudinal anchor and cross-mission predictive validation
have not yet been established. No discovery or validated forecast is reported.
\end{abstract}
\section{Introduction}
% Populate after the calibrated experiment and literature audit are extended.
\section{Radiation environments and missions}
% Separate LEO/L2 exposure, architecture, temperature and clocking in the future model.
\section{Data}
The provenance contract separates archival observations, genuinely published
calibration measurements and synthetic truth. MAST observation-query snapshots
define a deterministic paired-dark ACS/WFC pilot. Original files are retained in
a content-addressed local cache; source identifiers, retrieval times, software
revisions, processing state, usage notes and cryptographic hashes accompany them.
The release cohort is defined by the versioned machine-readable RAW product plan,
not by a newly evaluated archive search.

NASA SPDF OMNI hourly files have been retrieved as an ingestion pilot. The selected
years do not form a continuous exposure history. OMNI particle fluxes are
near-Earth environment measurements and are not detector dose. Missing entries
must remain missing. Granular public Gaia CTI and Euclid trap-pumping time series
have not been obtained. Public VIS imaging products and published calibration
constraints are different evidence classes \citep{pagani2026,skottfelt2024,q1vis}.
\section{Detector observables}
\subsection{Paired RAW-dark pilot}
Only individual full-frame dark exposures meeting the frozen time and exposure
criteria are selected. Array geometry is checked against the image headers and
each chip is oriented away from its parallel readout register. Because calibrated
gain entries are unpopulated in RAW headers, the pilot retains native digital
units and commanded gain as an explicit stratum.

The background is estimated from symmetric same-row side pixels. Persistent,
isolated peaks are required in both exposures. The observable is a downstream
minus upstream trail sum divided by the local peak signal. Leading, serial and
offset-column blank statistics are retained. Signal and transfer-distance bins
are fixed in the versioned method contract. Column-cluster bootstrap intervals
describe population sampling uncertainty, excluding calibration systematics.

This observable does not measure a physical trap density. Pixel bias structure,
readout effects, gain calibration, detector temperature and selection effects
remain unresolved. A CALACS/reference-calibrated comparison is required before
physical interpretation. The approach is informed by established warm-pixel and
capture/release methods \citep{anderson2010,massey2026}; it does not reproduce their
full calibrations.
\subsection{Official calibration and operating-state comparison}
Copies of the same RAW exposures are processed by the official ACSCCD stage,
using a pinned HSTCAL environment and CRDS context. Bias and overscan calibration
and detector gain conversion are applied, with the resulting geometry, units and
calibration switches checked explicitly. Dark, CTI and flash corrections are
omitted. The intermediates retain deferred charge and illumination effects and
are not fully calibrated science products. Source, reference, executable, log
and output hashes accompany each derived image.

The comparison uses separately frozen electron signal bins and excludes flagged
neighbourhoods except for hot/warm-pixel flags. Its selected population therefore
differs from the RAW native-unit pilot. Differences between the two summaries do
not isolate a causal calibration effect. Matched support files preserve both WFC
temperature channels without assuming an active-sensor mapping. Exposure headers
record post-flash duration and status, commanded gain and shutter position.
Background, temperature, electronics era and between-exposure variation remain
necessary nuisance variables. The initial serial statistic measures signed
array-axis asymmetry; opposite amplifier directions can cancel serial trails.
\subsection{Analytic injection validation}
Independent images with known exponential trails test extraction, sign convention,
blank controls and rejection of transient single-frame peaks. The injection is
labelled synthetic and is not a trap-occupancy simulation. It cannot establish
physical parameter recovery for a future latent-state model.
\section{Model}
% Not implemented: no observational posterior parameters to report.
\section{Validation}
% Final temporal cohort and holdout must be frozen before predictive evaluation.
\section{HST results}
% Intentionally unpopulated: RAW pilot has not passed the calibrated-anchor gate.
\section{Cross-mission results}
% No quantitative Gaia/Euclid transfer result exists.
\section{Science-impact propagation}
% No validated detector-to-science-bias experiment exists.
\section{Limitations}
% Expand after the calibrated experiment; current limitations are in the method above.
\section{Discussion}
\section{Conclusions}
\appendix
\section{Reproduction}
The repository includes exact commands, frozen manifests and a validation contract.
Any future numerical manuscript result must be generated from a versioned results
file. The present document contains no empirical result claims.
\bibliographystyle{plainnat}
\bibliography{references}
\end{document}
